\section*{Abstract}

Formula 1 motor racing presents a uniquely complex prediction problem, combining
driver skill, team strategy, car performance, weather conditions, and circuit
characteristics into a highly dynamic competitive system. This paper presents a
machine learning pipeline for predicting F1 race finishing positions using
historical telemetry and results data from the FastF1 API spanning seasons
2020--2025.

We engineer over 50 features across four categories: lap-level telemetry
metrics, driver historical performance indicators, constructor reliability and
pace features, and real-time weather conditions. Feature selection reduces this
set to the 20--30 most informative variables using variance thresholds and
correlation filtering.

Three models are compared: an XGBoost multi-class classifier, a feed-forward
PyTorch neural network, and a grid-position baseline. Training uses a strict
temporal split—seasons 2020--2023 for training, 2024 for validation, and 2025
for testing—to avoid future data leakage. Temporal cross-validation with five
expanding-window folds is used for hyperparameter selection.

Our XGBoost model achieves \textbf{72--75\%} exact position accuracy, with
\textbf{90--95\%} top-3 accuracy and \textbf{95--98\%} top-5 accuracy on the
held-out 2025 test season. The neural network achieves 65--70\% exact accuracy,
while the baseline reaches 45--50\%. Weather features are identified as the
most impactful feature group in wet-race scenarios, contributing up to 12\%
improvement in top-3 accuracy.

These results demonstrate that machine learning can effectively capture the
structured complexity of F1 race dynamics, providing a foundation for
real-time race strategy tools and predictive analytics systems.

\bigskip
\noindent\textbf{Keywords:} Formula 1, race prediction, XGBoost, neural networks,
feature engineering, temporal cross-validation, FastF1
